\chapter{Đọc thêm 1: Công nghệ khoa học dữ liệu trong kinh tế và tài chính}
\label{ch:M1L1-reading1}

% Nguồn: Consoli, S., Barbaglia, L., Manzan, S., Reforgiato Recupero, D., Saisana, M., Tiozzo Pezzoli, L.
% "Data Science Technologies in Economics and Finance: A Gentle Walk-In"
% trong: Data Science for Economics and Finance: Methodologies and Applications, Springer Nature, 2021.
% https://link.springer.com/book/10.1007/978-3-030-66891-4
% Bài viết gốc được cấp phép Creative Commons Attribution 4.0 International (CC BY 4.0),
% cho phép sử dụng, chia sẻ, chuyển thể, phân phối và sao chép dưới mọi hình thức, miễn là
% ghi công tác giả gốc và nguồn, có liên kết tới giấy phép CC và nêu rõ nếu có thay đổi.
% Bản dịch tiếng Việt dưới đây là bản dịch toàn văn của chương, phục vụ mục đích học tập.

\section*{Tóm tắt}

Chương này là phần giới thiệu về việc sử dụng các công nghệ khoa học dữ liệu trong lĩnh vực kinh tế và tài chính. Sự bùng nổ gần đây của công nghệ tính toán và thông tin trong thập kỷ qua đã tạo ra một lượng dữ liệu khổng lồ trong nhiều lĩnh vực khác nhau, thường được gọi là \emph{Dữ liệu lớn} (Big Data). Trong kinh tế và tài chính nói riêng, việc khai thác các nguồn dữ liệu này giúp nghiên cứu và hoạt động kinh doanh xích lại gần nhau hơn, vì dữ liệu được sinh ra trong hoạt động kinh tế thường nhật có thể được sử dụng để xây dựng các mô hình hiệu quả và mang tính cá nhân hóa. Trong bối cảnh đó, việc ứng dụng gần đây các công nghệ khoa học dữ liệu cho kinh tế và tài chính đem lại lợi ích cho cả giới nghiên cứu lẫn giới chuyên môn, giúp cải thiện khả năng dự báo (forecasting) và ước tính hiện tại (nowcasting) cho nhiều loại ứng dụng khác nhau. Chương này giới thiệu chủ đề thông qua các thách thức kỹ thuật nền tảng như quản lý và bảo vệ dữ liệu, mô hình hóa, tích hợp và diễn giải kết quả. Chương cũng phác thảo một số vấn đề phổ biến trong mô hình hóa kinh tế bằng công nghệ khoa học dữ liệu, đồng thời khảo sát các giải pháp quản lý và phân tích dữ liệu lớn liên quan, nhằm làm rõ động lực cho việc ứng dụng các phương pháp khoa học dữ liệu trong kinh tế và tài chính.

\section{Giới thiệu}

Những tiến bộ nhanh chóng của công nghệ thông tin và truyền thông trong hai thập kỷ qua đã tạo ra sự tăng trưởng bùng nổ về lượng thông tin được thu thập, dẫn tới kỷ nguyên mới của dữ liệu lớn~[31]. Theo~[26], mỗi ngày có khoảng ba tỷ byte dữ liệu được tạo ra từ các cảm biến, thiết bị di động, giao dịch trực tuyến và mạng xã hội, với 90\% lượng dữ liệu trên thế giới được tạo ra chỉ trong ba năm gần đây. Những thách thức về lưu trữ, tổ chức và thấu hiểu một khối lượng thông tin khổng lồ như vậy đã thúc đẩy sự phát triển của các công nghệ mới trong nhiều lĩnh vực khác nhau như thống kê, học máy và khai phá dữ liệu, đồng thời tương tác với các lĩnh vực kỹ thuật và trí tuệ nhân tạo (AI), trong số nhiều lĩnh vực khác. Nỗ lực to lớn này đã khai sinh ra một lĩnh vực liên ngành mới gọi là "Khoa học dữ liệu" (Data Science), với các nguyên lý và kỹ thuật nhằm tự động trích xuất thông tin và tri thức hữu ích tiềm tàng từ dữ liệu. Mặc dù các công nghệ khoa học dữ liệu đã được ứng dụng thành công trong nhiều lĩnh vực khác nhau (ví dụ như y tế~[15], bảo trì dự đoán~[16], và quản lý chuỗi cung ứng~[39], trong số nhiều lĩnh vực khác), tiềm năng của chúng vẫn chưa được khai thác nhiều trong kinh tế và tài chính. Trong bối cảnh này, việc xây dựng các mô hình dự báo và ước tính hiện tại hiệu quả là điều thiết yếu để thiết kế các chính sách tiền tệ và tài khóa phù hợp, và độ chính xác của các mô hình này đặc biệt quan trọng trong những giai đoạn kinh tế biến động. Việc theo dõi trạng thái hiện tại và tương lai của nền kinh tế có tầm quan trọng nền tảng đối với các chính phủ, tổ chức quốc tế và ngân hàng trung ương trên toàn thế giới. Các nhà hoạch định chính sách cần có thông tin kinh tế vĩ mô sẵn sàng để thiết kế những chính sách hiệu quả nhằm thúc đẩy tăng trưởng kinh tế và duy trì phúc lợi xã hội. Tuy nhiên, các chỉ số kinh tế then chốt mà họ dựa vào trong quá trình ra quyết định thường được công bố với tần suất thấp và có độ trễ đáng kể — ví dụ, khoảng 45 ngày đối với Tổng sản phẩm quốc nội (GDP) tại châu Âu — và thường phải chịu những lần điều chỉnh có thể rất lớn. Thật vậy, với một tập hợp thông tin không đầy đủ như vậy, các nhà kinh tế chỉ có thể ước lượng một cách tương đối điều kiện kinh tế thực tế, tương lai, và thậm chí cả quá khứ gần, khiến việc ước tính hiện tại và dự báo nền kinh tế trở thành những nhiệm vụ vô cùng thách thức. Thêm vào đó, trong một thế giới toàn cầu hóa và liên kết chặt chẽ, các cú sốc và biến động bắt nguồn từ một nền kinh tế có thể nhanh chóng lan sang các nền kinh tế khác, ảnh hưởng đến năng suất, tạo việc làm và phúc lợi ở nhiều khu vực địa lý khác nhau. Tóm lại, các nhà hoạch định chính sách phải đối mặt với một bài toán kép: tính kịp thời trong việc đánh giá nền kinh tế, cũng như việc đánh giá nhanh chóng tác động của các cú sốc bên ngoài.

Các mô hình dự báo truyền thống áp dụng phương pháp tần suất hỗn hợp, kết hợp thông tin từ các chỉ số kinh tế và tài chính tần suất cao (ví dụ như sản lượng công nghiệp hoặc giá cổ phiếu) cũng như các khảo sát kinh tế với biến số tần suất thấp cần dự báo, chẳng hạn như GDP~[28]. Một cách tiếp cận khác có thể là các mô hình nhân tố động (dynamic factor models), vốn tóm lược một lượng lớn thông tin thành một vài nhân tố và xử lý dữ liệu bị thiếu bằng kỹ thuật lọc Kalman trong quá trình ước lượng. Những cách tiếp cận này cho phép sử dụng hàm phản ứng xung (impulse-response) để đánh giá phản ứng của nền kinh tế trước các cú sốc bên ngoài, cung cấp những định hướng chung cho các nhà hoạch định chính sách để xây dựng chính sách hiện tại và tương lai, có tính đến đầy đủ thông tin đến từ bên ngoài. Tuy nhiên, các phương pháp truyền thống này có hai hạn chế chính. Thứ nhất, chúng không thể xử lý trực tiếp một lượng lớn dữ liệu phi cấu trúc vì được thiết kế riêng cho các nguồn dữ liệu có cấu trúc. Thứ hai, ngay cả khi các mô hình cổ điển này được bổ sung thêm các biến dự báo mới thu được từ các tập dữ liệu lớn thay thế, mối quan hệ giữa các biến vẫn được giả định là tuyến tính, điều này không đúng với phần lớn các trường hợp trong thực tế~[21, 1].

Các công nghệ khoa học dữ liệu cho phép các nhà kinh tế giải quyết tất cả những vấn đề này. Một mặt, các nguồn dữ liệu lớn mới có thể tích hợp và bổ sung thông tin do các biến số tổng hợp công khai được sản xuất bởi các cơ quan thống kê quốc gia và quốc tế cung cấp. Mặt khác, các thuật toán học máy có thể trích xuất những hiểu biết mới từ những thông tin phi cấu trúc đó và tính đến một cách phù hợp các động lực phi tuyến giữa các biến kinh tế và tài chính. Về mặt dữ liệu lớn, mức độ chi tiết cao hơn được thể hiện trong các nguồn dữ liệu mới có sẵn tạo ra tiềm năng mạnh mẽ để khám phá những mối quan hệ kinh tế thường không rõ ràng khi các biến được tổng hợp trên nhiều sản phẩm, cá nhân hoặc giai đoạn thời gian. Một số ví dụ về các nguồn dữ liệu lớn mới có thể hữu ích cho việc dự báo và ước tính hiện tại kinh tế bao gồm: dữ liệu giá bán lẻ từ máy quét (scanner price data), giao dịch thẻ tín dụng/ghi nợ, đồng hồ đo năng lượng thông minh, cảm biến giao thông thông minh, ảnh vệ tinh, tin tức thời gian thực và dữ liệu mạng xã hội. Dữ liệu giá từ máy quét, giao dịch thẻ và đồng hồ thông minh cung cấp thông tin về người tiêu dùng, từ đó mang lại khả năng hiểu rõ hơn hành vi thực tế của các biến số vĩ mô như GDP hay các thành phần của lạm phát. Ảnh vệ tinh và cảm biến giao thông có thể được dùng để theo dõi phương tiện thương mại, tàu thuyền và hoạt động của nhà máy, khiến chúng trở thành nguồn dữ liệu tiềm năng để ước tính sản lượng công nghiệp. Tin tức thời gian thực và mạng xã hội có thể được sử dụng để đại diện cho tâm trạng của các chủ thể kinh tế và tài chính, và có thể được xem như một thước đo về nhận thức đối với trạng thái thực tế của nền kinh tế.

Bên cạnh dữ liệu mới, các phương pháp thay thế như thuật toán học máy có thể giúp các nhà kinh tế mô hình hóa các hệ thống động phức tạp và liên kết chặt chẽ với nhau. Chúng có khả năng nắm bắt tri thức ẩn ngay cả khi số lượng đặc trưng (feature) đang được phân tích lớn hơn số lượng quan sát sẵn có, điều thường xảy ra trong môi trường kinh tế. Khác với các kỹ thuật chuỗi thời gian truyền thống, các phương pháp học máy không có giả định "tiên nghiệm" nào về quá trình ngẫu nhiên nằm bên dưới trạng thái của nền kinh tế. Chẳng hạn, học sâu (deep learning)~[29], một phương pháp khoa học dữ liệu rất phổ biến hiện nay, hữu ích trong việc mô hình hóa dữ liệu có tính phi tuyến cao vì bậc phi tuyến được suy ra hoặc học trực tiếp từ dữ liệu chứ không phải được giả định sẵn như trong nhiều mô hình kinh tế lượng truyền thống. Các mô hình khoa học dữ liệu có khả năng khám phá những mối quan hệ phức tạp, điều này có thể hữu ích để dự báo và ước tính hiện tại nền kinh tế trong thời kỳ bình thường, nhưng cũng để phát hiện sớm các dấu hiệu bất ổn trên thị trường trước khi xảy ra khủng hoảng tài chính.

Mặc dù những phương pháp này có thể mang lại các dự báo chính xác, việc hiểu được ý nghĩa kinh tế đằng sau những kết quả đầy hứa hẹn đó lại là một nhiệm vụ khó khăn. Bản chất của các phương pháp này là "hộp đen" (black box), được phát triển với mục tiêu duy nhất là tối đa hóa hiệu năng dự báo. Toàn bộ lĩnh vực khoa học dữ liệu được đánh giá dựa trên các thử nghiệm ngoài mẫu (out-of-sample), xem một mô hình được huấn luyện trên một tập dữ liệu sẽ dự báo dữ liệu mới tốt đến mức nào. Ngược lại, các nhà kinh tế cần biết các mô hình có thể tác động ra sao trong thế giới thực, và họ thường không chỉ tập trung vào dự báo mà còn quan tâm đến suy luận mô hình (model inference), tức là hiểu các tham số của mô hình (ví dụ, kiểm định các hệ số riêng lẻ trong một hồi quy). Các nhà hoạch định chính sách cần có cơ sở để bảo vệ quyết định của mình và đưa ra một tập hợp các lời giải thích khả dĩ cho một hành động đã thực hiện; do đó, họ quan tâm đến hàm ý kinh tế liên quan đến các dự báo của mô hình. Các hàm phản ứng xung là công cụ quen thuộc để đánh giá tác động của một cú sốc lên một biến đối với kết quả quan tâm, nhưng các thuật toán học máy không hỗ trợ chức năng này. Điều này có thể cản trở, chẳng hạn, việc đánh giá các chính sách ổn định nhằm bảo vệ nhu cầu nội địa khi một cú sốc bên ngoài tác động vào nền kinh tế. Để lấp đầy khoảng trống này, cộng đồng khoa học dữ liệu gần đây đã nỗ lực tăng cường tính minh bạch của các mô hình học máy trong tài liệu về \emph{AI có thể diễn giải} (interpretable AI)~[22]. Các ứng dụng học máy trong kinh tế và tài chính giờ đây có thể hưởng lợi từ những công cụ mới như biểu đồ Phụ thuộc riêng phần (Partial Dependence plots) hay giá trị Shapley, cho phép các nhà hoạch định chính sách đánh giá tác động biên của các biến mô hình lên kết quả dự báo. Tóm lại, khoa học dữ liệu có thể tăng cường các mô hình dự báo kinh tế bằng cách:

\begin{itemize}
\item Tích hợp và bổ sung các chỉ số thống kê chính thức bằng cách sử dụng các nguồn dữ liệu lớn phi cấu trúc, thời gian thực mới
\item Đánh giá điều kiện kinh tế và tài chính hiện tại và tương lai bằng cách cho phép các mối quan hệ phi tuyến phức tạp giữa các biến dự báo
\item Tối đa hóa doanh thu của giao dịch thuật toán (algorithmic trading), một nhiệm vụ hoàn toàn dựa trên dữ liệu
\item Hỗ trợ ra quyết định một cách phù hợp bằng cách làm cho kết quả đầu ra của các thuật toán học máy trở nên dễ hiểu
\end{itemize}

Chương này nhấn mạnh rằng khoa học dữ liệu có tiềm năng giải phóng những nút thắt năng suất to lớn và cải thiện triệt để chất lượng cũng như khả năng tiếp cận của các mô hình dự báo kinh tế, đồng thời thảo luận về những thách thức và các bước cần được cân nhắc để đảm bảo việc ứng dụng rộng rãi và chuyên sâu.

\section{Các thách thức kỹ thuật}

Trong những năm gần đây, những tiến bộ công nghệ đã làm gia tăng đáng kể số lượng thiết bị tạo ra thông tin về hoạt động của con người và kinh tế (ví dụ: cảm biến, giám sát, thiết bị IoT, mạng xã hội). Những nguồn dữ liệu mới này cung cấp một lượng thông tin phong phú, tần suất cao và đa dạng, từ đó trạng thái của nền kinh tế có thể được ước lượng với độ chính xác và tính kịp thời. Việc thu thập và phân tích những loại dữ liệu này là một nhiệm vụ đầy thách thức do kích thước và sự đa dạng của chúng. Tuy nhiên, nếu được khai thác đúng cách, những nguồn dữ liệu mới này có thể mang lại sức mạnh dự báo bổ sung so với các biến hồi quy tiêu chuẩn được sử dụng trong phân tích kinh tế và tài chính truyền thống.

Khi kích thước và sự đa dạng của dữ liệu tăng lên, nhu cầu về những cỗ máy mạnh mẽ hơn và các thuật toán hiệu quả hơn trở nên rõ ràng hơn. Việc phân tích những loại dữ liệu này có thể đòi hỏi tính toán rất cao và đã tạo ra nhu cầu ngày càng tăng đối với phần cứng và môi trường điện toán hiệu quả. Chẳng hạn, các Bộ xử lý đồ họa (GPU) và hệ thống điện toán đám mây trong những năm gần đây đã trở nên phải chăng hơn và được sử dụng bởi đông đảo người dùng hơn. GPU có kiến trúc song song dữ liệu cao, có thể được lập trình bằng các khung công cụ như CUDA\footnote{NVIDIA CUDA: \url{https://developer.nvidia.com/cuda-zone}.} và OpenCL\footnote{OpenCL: \url{https://www.khronos.org/opencl/}.}. Chúng bao gồm một số lõi (core), mỗi lõi có một số đơn vị chức năng. Một hoặc nhiều đơn vị chức năng này (gọi là \emph{bộ xử lý luồng}, thread processor) xử lý từng luồng thực thi. Tất cả các bộ xử lý luồng trong một lõi của GPU thực hiện cùng một tập lệnh, vì chúng dùng chung một đơn vị điều khiển. Điện toán đám mây đại diện cho việc phân phối các dịch vụ như máy chủ, cơ sở dữ liệu và phần mềm thông qua Internet. Về cơ bản, một nhà cung cấp cấp cho người dùng quyền truy cập theo nhu cầu vào các dịch vụ lưu trữ, xử lý và truyền tải dữ liệu. Một số ví dụ về giải pháp điện toán đám mây là Google Cloud Platform\footnote{Google Cloud: \url{https://cloud.google.com/}.}, Microsoft Azure\footnote{Microsoft Azure: \url{https://azure.microsoft.com/en-us/}.} và Amazon Web Services (AWS)\footnote{Amazon Web Services (AWS): \url{https://aws.amazon.com/}.}.

Sức mạnh tính toán đủ lớn là điều kiện cần để phân tích các nguồn dữ liệu lớn mới; tuy nhiên, điều đó chưa đủ nếu dữ liệu không được lưu trữ, chuyển đổi và kết hợp một cách phù hợp. Hiện nay, các tập dữ liệu kinh tế và tài chính vẫn được lưu trữ trong các "kho riêng lẻ" (silo), và các nhà nghiên cứu cũng như người thực hành thường gặp khó khăn khi kết hợp chúng giữa nhiều nhà cung cấp, các tổ chức kinh tế khác, và thậm chí cả dữ liệu do người tiêu dùng tạo ra. Những tập dữ liệu kinh tế khác biệt này có thể khác nhau về mức độ chi tiết, chất lượng và loại dữ liệu, chẳng hạn từ văn bản tự do, hình ảnh, dữ liệu cảm biến (streaming) cho đến các tập dữ liệu có cấu trúc; việc tích hợp chúng đặt ra những thách thức lớn về pháp lý, kinh doanh và kỹ thuật. Dữ liệu lớn và công nghệ khoa học dữ liệu hướng đến việc giải quyết hiệu quả những thách thức này.

Thuật ngữ "dữ liệu lớn" (big data) có nguồn gốc từ ngành kỹ thuật máy tính. Mặc dù có nhiều định nghĩa về dữ liệu lớn trong tài liệu~[31, 43], chúng ta có thể hiểu một cách trực quan rằng đó là dữ liệu lớn đến mức không thể nạp vào bộ nhớ hoặc thậm chí không thể lưu trữ trên một máy đơn lẻ. Ngoài \emph{khối lượng} (volume) lớn, còn có các chiều khác đặc trưng cho dữ liệu lớn, đó là \emph{sự đa dạng} (variety, xử lý nhiều loại, nguồn và định dạng khác nhau), \emph{độ xác thực} (veracity, liên quan đến chất lượng và tính hợp lệ của dữ liệu), và \emph{tốc độ} (velocity, khả năng sẵn có của dữ liệu theo thời gian thực). Ngoài bốn đặc điểm dữ liệu lớn nói trên, chúng ta cũng cần xem xét các vấn đề liên quan như độ tin cậy của dữ liệu, bảo vệ dữ liệu và quyền riêng tư dữ liệu. Trong chương này, chúng ta sẽ khám phá những thách thức chính do việc khai thác các nguồn dữ liệu mới và thay thế đặt ra, cùng với những giải pháp tương ứng mà cộng đồng khoa học dữ liệu đã xây dựng.

\subsection{Quản lý và bảo vệ dữ liệu (Stewardship and Protection)}

Khả năng tiếp cận là một điều kiện quan trọng để khai thác hiệu quả các nguồn dữ liệu mới cho phân tích kinh tế và tài chính. Tuy nhiên, trong thực tế, khả năng tiếp cận thường bị hạn chế nhằm bảo vệ thông tin nhạy cảm. Việc tìm ra một sự cân bằng hợp lý giữa khả năng tiếp cận và bảo vệ dữ liệu thường được gọi là \emph{quản lý dữ liệu} (data stewardship), một khái niệm bao gồm từ việc thu thập, chú thích và lưu trữ thông tin đúng cách cho đến việc "chăm sóc lâu dài" dữ liệu, coi dữ liệu như những tài sản số có giá trị có thể được tái sử dụng trong các ứng dụng tương lai và kết hợp với dữ liệu mới~[42]. Các tổ chức như World Wide Web Consortium (W3C)\footnote{World Wide Web Consortium (W3C): \url{https://www.w3.org/}.} đã nỗ lực xây dựng các hướng dẫn về khả năng tương tác (interoperability) trong phạm vi các tập dữ liệu mở có sẵn ở nhiều lĩnh vực khác nhau, nhằm đảm bảo rằng dữ liệu tuân theo nguyên tắc FAIR (\emph{F}indable — có thể tìm thấy, \emph{A}ccessible — có thể truy cập, \emph{I}nteroperable — có thể tương tác, và \emph{R}eusable — có thể tái sử dụng).

Bảo vệ dữ liệu là một khía cạnh then chốt cần được xem xét khi làm việc với dữ liệu kinh tế và tài chính. Độ tin cậy là mối quan tâm chính của cá nhân và tổ chức khi đối diện với việc sử dụng dữ liệu liên quan đến tài chính của họ: điều quan trọng là dữ liệu đó phải được lưu trữ trong các cơ sở dữ liệu an toàn và tôn trọng quyền riêng tư. Hiện nay, có nhiều phương pháp bảo vệ quyền riêng tư khác nhau để phân tích một nguồn dữ liệu cụ thể hoặc để kết nối các cơ sở dữ liệu khác nhau giữa các lĩnh vực hoặc kho lưu trữ. Tuy vậy, vẫn còn nhiều thách thức và rủi ro cần được giải quyết nhằm kết hợp các cơ sở dữ liệu riêng tư bằng các phương pháp ẩn danh hóa và giả ẩn danh hóa mới, đảm bảo quyền riêng tư. Các kỹ thuật phân tích dữ liệu cần được điều chỉnh để làm việc với dữ liệu được mã hóa hoặc phân tán. Sự hợp tác chặt chẽ giữa các chuyên gia trong lĩnh vực và các nhà phân tích dữ liệu trong suốt toàn bộ chuỗi khoa học dữ liệu có tầm quan trọng đặc biệt.

Dữ liệu ở cấp độ cá nhân về hiệu suất tín dụng là một ví dụ rõ ràng về dữ liệu nhạy cảm có thể rất hữu ích trong phân tích kinh tế và tài chính, nhưng việc tiếp cận thường bị hạn chế vì lý do bảo vệ dữ liệu. Việc khai thác đúng cách những dữ liệu này có thể mang lại những cải thiện lớn ở nhiều khía cạnh: các tổ chức tài chính có thể hưởng lợi từ những mô hình rủi ro tín dụng tốt hơn, giúp xác định chính xác hơn những người vay có rủi ro cao và giảm thiểu tổn thất tiềm tàng liên quan đến vỡ nợ; người tiêu dùng có thể dễ dàng tiếp cận tín dụng hơn nhờ việc phân bổ nguồn lực hiệu quả cho những người vay đáng tin cậy; và các chính phủ cùng ngân hàng trung ương có thể giám sát theo thời gian thực tình trạng của nền kinh tế bằng cách kiểm tra sức khỏe của thị trường tín dụng. Có rất nhiều tập dữ liệu chứa thông tin cấp độ cá nhân đã được ẩn danh hóa sẵn có trực tuyến. Chẳng hạn, dữ liệu thế chấp (mortgage) tại Hoa Kỳ được cung cấp bởi Hiệp hội Thế chấp Quốc gia Liên bang (Fannie Mae)\footnote{Federal National Mortgage Association (Fannie Mae): \url{https://www.fanniemae.com}.} và Tập đoàn Thế chấp Nhà ở Liên bang (Freddie Mac)\footnote{Federal Home Loan Mortgage Corporation (Freddie Mac): \url{http://www.freddiemac.com}.}: các tổ chức này công bố thông tin ở cấp độ từng khoản vay cho hàng triệu khoản thế chấp cá nhân, với nhiều đặc trưng liên quan, ví dụ như tình trạng trả nợ, đặc điểm chính của người vay, và địa điểm cấp khoản vay (xem~[2, 35] để biết hai ví dụ về phân tích cấp độ thế chấp tại Hoa Kỳ). Một mức độ chi tiết tương tự cũng được tìm thấy tại European Datawarehouse\footnote{European Datawarehouse: \url{https://www.eurodw.eu/}.}, nơi cung cấp dữ liệu cấp độ khoản vay của các tài sản châu Âu liên quan đến thế chấp nhà ở, thẻ tín dụng, cho thuê ô tô và tài chính tiêu dùng (xem~[20, 40] để biết hai ví dụ về phân tích kinh tế dựa trên loại dữ liệu này).

\subsection{Số lượng dữ liệu và dữ liệu chuẩn (Ground Truth)}

Dữ liệu kinh tế và tài chính đang tăng trưởng với tốc độ chưa từng thấy trong quá khứ~[33]. Ngày nay, các tổ chức đang thu thập một khối lượng lớn dữ liệu từ cả nguồn sở hữu riêng lẫn nguồn công khai, chẳng hạn như mạng xã hội và dữ liệu mở, và cuối cùng sử dụng chúng cho phân tích kinh tế và tài chính. Khối lượng và tốc độ dữ liệu ngày càng tăng đặt ra những thách thức kỹ thuật mới mà các nhà nghiên cứu và phân tích có thể giải quyết bằng cách tận dụng khoa học dữ liệu. Một kịch bản khoa học dữ liệu tổng quát bao gồm một chuỗi các quan sát, thường được gọi là các \emph{thực thể} (instance), mỗi thực thể được đặc trưng bởi giá trị hiện thực hóa của một nhóm biến, thường gọi là các \emph{thuộc tính} (attribute), có thể có dạng như một chuỗi văn bản, một mã chữ-số, một ngày tháng, một thời điểm, hoặc một con số. Khối lượng dữ liệu đang bùng nổ theo nhiều hướng khác nhau: ngày càng có nhiều tập dữ liệu sẵn có hơn, mỗi tập có số lượng thực thể ngày càng tăng; những tiến bộ công nghệ cho phép thu thập thông tin trên một số lượng lớn đặc trưng, kể cả dưới dạng hình ảnh và video.

Các nhà khoa học dữ liệu thường phân biệt hai loại dữ liệu: dữ liệu không có nhãn (unlabeled) và dữ liệu có nhãn (labeled)~[15]. Với một thuộc tính quan tâm (nhãn) cho trước, dữ liệu không có nhãn là dữ liệu không đi kèm giá trị quan sát của nhãn đó, và chúng được sử dụng trong các bài toán học không giám sát (unsupervised learning), nơi mục tiêu là trích xuất tối đa thông tin có sẵn từ chính dữ liệu, chẳng hạn như các bài toán phân cụm (clustering) và luật kết hợp (association rules)~[15]. Đối với loại dữ liệu thứ hai, mỗi thực thể dữ liệu đi kèm với một nhãn có thể được sử dụng trong một nhiệm vụ học có giám sát (supervised learning): người ta có thể sử dụng thông tin sẵn có trong tập dữ liệu để dự đoán giá trị của thuộc tính quan tâm chưa được quan sát. Nếu thuộc tính quan tâm là biến định tính (categorical), nhiệm vụ đó gọi là phân loại (classification); còn nếu là biến định lượng (numerical), nhiệm vụ đó gọi là hồi quy (regression)~[15]. Các công nghệ đột phá như học sâu đòi hỏi một lượng lớn dữ liệu có nhãn để huấn luyện, nghĩa là dữ liệu cần đi kèm với các chú thích, thường được gọi là \emph{dữ liệu chuẩn} (ground truth)~[15].

Trong lĩnh vực tài chính, ví dụ, nhiều công trình về học không giám sát và học có giám sát đã được khai thác trong tài liệu về phát hiện gian lận~[3, 11], với mục tiêu xác định xem một giao dịch tài chính nhất định có xảy ra gian lận tiềm tàng hay không. Trong lĩnh vực này, bộ dữ liệu nổi tiếng về Phát hiện gian lận thẻ tín dụng (Credit Card Fraud Detection)\footnote{\url{https://www.kaggle.com/mlg-ulb/creditcardfraud}.} thường được sử dụng để so sánh hiệu năng của các thuật toán khác nhau trong việc nhận diện hành vi gian lận (ví dụ~[17, 32]). Bộ dữ liệu này chứa 284.807 giao dịch của các chủ thẻ châu Âu được thực hiện trong 2 ngày của năm 2013, trong đó chỉ có 492 giao dịch được đánh dấu là gian lận, tức 0,17\% tổng số. Số lượng trường hợp dương tính nhỏ này cần được phân chia một cách nhất quán vào tập huấn luyện và tập kiểm tra thông qua lấy mẫu phân tầng (stratified sampling), sao cho cả hai tập đều chứa một số giao dịch gian lận nhằm đảm bảo so sánh công bằng về hiệu năng dự báo ngoài mẫu. Do khối lượng dữ liệu ngày càng tăng, việc làm việc với những tập dữ liệu mất cân bằng cao như vậy ngày càng trở nên phổ biến, trong đó số lượng trường hợp dương tính chỉ chiếm một phần nhỏ của toàn bộ tập dữ liệu: trong những trường hợp này, phân tích kinh tế lượng tiêu chuẩn có thể mang lại kết quả kém, và việc xem xét các kỹ thuật cân bằng lại như lấy mẫu dưới (undersampling), lấy mẫu trên (oversampling), hoặc kết hợp cả hai, có thể hữu ích nhằm cải thiện độ chính xác phân loại~[15, 36].

\subsection{Chất lượng dữ liệu và nguồn gốc dữ liệu (Provenance)}

Chất lượng dữ liệu nhìn chung đề cập đến việc dữ liệu nhận được có phù hợp với mục đích sử dụng và phân tích dự kiến hay không. Cơ sở để đánh giá chất lượng dữ liệu được cung cấp là phải có một phần siêu dữ liệu (metadata) được cập nhật, trong đó có mô tả đầy đủ cho từng đặc trưng trong phân tích. Cần nhấn mạnh rằng phần lớn công việc của nhà khoa học dữ liệu nằm ở việc kiểm tra xem các bản ghi dữ liệu có thực sự tương ứng với các mô tả trong siêu dữ liệu hay không. Lỗi của con người và dữ liệu không nhất quán hoặc bị thiên lệch có thể tạo ra sự khác biệt so với những gì bên nhận dữ liệu kỳ vọng ban đầu. Chẳng hạn, đối với European Datawarehouse được trình bày ở Mục~2.1: dữ liệu cấp độ khoản vay được từng tổ chức tài chính báo cáo, tập hợp trên một nền tảng tập trung và công bố theo một cấu trúc dữ liệu chung. Các tổ chức tài chính được hướng dẫn đầy đủ về cách cung cấp dữ liệu; tuy nhiên, nhiều loại lỗi khác nhau vẫn có thể xảy ra. Ví dụ, tỷ lệ lãi suất có thể được báo cáo dưới dạng phân số thay vì phần trăm, và các khoản vay có thể được đánh dấu là vỡ nợ theo một định nghĩa thay đổi theo thời gian và/hoặc theo luật pháp của từng quốc gia.

Đi xa hơn các kiểm tra chất lượng dữ liệu tiêu chuẩn, \emph{nguồn gốc dữ liệu} (data provenance) hướng đến việc thu thập thông tin về toàn bộ quá trình sinh ra dữ liệu, chẳng hạn như phần mềm được sử dụng, các bước thực nghiệm đã thực hiện trong việc thu thập dữ liệu, hoặc bất kỳ chi tiết nào về các thao tác trước đó đã được thực hiện trên dữ liệu thô đầu vào. Việc theo dõi thông tin này cho phép bên nhận dữ liệu hiểu được nguồn gốc của dữ liệu, tức là dữ liệu được thu thập như thế nào, trong điều kiện nào, cũng như được xử lý và chuyển đổi ra sao trước khi được lưu trữ. Hơn nữa, nếu bên cung cấp dữ liệu thay đổi bất kỳ khía cạnh nào được xem xét trong nguồn gốc dữ liệu (ví dụ, cập nhật phần mềm), bên nhận dữ liệu có thể phát hiện sớm một sự thay đổi cấu trúc trong chất lượng dữ liệu, từ đó ngăn ngừa việc sử dụng và phân tích sai lệch tiềm tàng. Điều này quan trọng không chỉ đối với khả năng tái lập (reproducibility) của phân tích, mà còn đối với việc hiểu độ tin cậy của dữ liệu, vốn có thể ảnh hưởng đến kết quả trong nghiên cứu kinh tế. Khi độ phức tạp của các thao tác ngày càng tăng, cùng với các phương pháp mới được phát triển khá nhanh, việc ghi lại và hiểu nguồn gốc của dữ liệu trở nên then chốt, vì điều này có thể ảnh hưởng đáng kể đến kết luận của phân tích. Để có một bài tổng quan gần đây về tương lai của nguồn gốc dữ liệu, xin xem, trong số nhiều tài liệu khác,~[10].

\subsection{Tích hợp và chia sẻ dữ liệu}

Khoa học dữ liệu làm việc với dữ liệu có cấu trúc và phi cấu trúc được sinh ra từ nhiều nguồn khác nhau và ở nhiều định dạng khác nhau, hướng đến việc tích hợp chúng vào các kho dữ liệu lớn hoặc Kho dữ liệu (Data Warehouse)~[43]. Có một số lượng lớn các thao tác ETL (Trích xuất, Chuyển đổi và Nạp — Extraction, Transformation, and Loading) được chuẩn hóa, giúp xác định và tổ chức lại tính không đồng nhất về cấu trúc, cú pháp và ngữ nghĩa giữa các nguồn dữ liệu khác nhau~[31]. Tính không đồng nhất về cấu trúc đề cập đến các mô hình dữ liệu và lược đồ khác nhau, đòi hỏi tích hợp ở cấp độ lược đồ. Tính không đồng nhất về cú pháp xuất hiện dưới dạng các giao diện truy cập dữ liệu khác nhau, cần được thống nhất lại. Tính không đồng nhất về ngữ nghĩa nằm ở sự khác biệt trong cách diễn giải giá trị dữ liệu, và có thể được khắc phục bằng cách sử dụng các công nghệ ngữ nghĩa, như cơ sở tri thức dựa trên đồ thị (graph-based knowledge base) và bản thể miền (domain ontology)~[8], giúp ánh xạ các khái niệm và định nghĩa với nguồn dữ liệu, từ đó tạo điều kiện cho sự hợp tác, chia sẻ, mô hình hóa và tái sử dụng giữa các ứng dụng~[7].

Một quá trình tích hợp cuối cùng dẫn đến việc hợp nhất các nguồn và tập dữ liệu trùng lặp. Việc tích hợp và liên kết dữ liệu có thể được tăng cường hơn nữa bằng cách khai thác đúng cách các thuật toán trích xuất thông tin, phương pháp học máy và các công nghệ Web ngữ nghĩa (Semantic Web) cho phép diễn giải thông tin dựa trên ngữ cảnh~[26]. Chẳng hạn, các tác giả trong~[12] đề xuất một phương pháp ngữ nghĩa để tạo ra các từ điển đặc thù theo ngành từ các tài liệu tin tức được thu thập trong bộ dữ liệu Dow Jones DNA\footnote{Dow Jones DNA: \url{https://www.dowjones.com/dna/}.}, với mục tiêu nắm bắt một cách động, theo từng ngày, mối tương quan giữa các từ được sử dụng trong những tài liệu này và biến động giá cổ phiếu của các ngành trong chỉ số Standard \& Poor's 500. Một ví dụ khác được trình bày trong công trình~[37], sử dụng thông tin trích xuất từ \emph{Wall Street Journal} để cho thấy rằng mức độ bi quan cao trong tin tức là những yếu tố dự báo quan trọng cho sự hội tụ của giá cổ phiếu về giá trị nền tảng của chúng.

Trong lĩnh vực kinh tế vĩ mô,~[24] đã nghiên cứu nội dung thông tin trong các tuyên bố của Cục Dự trữ Liên bang (Federal Reserve) và định hướng mà những tuyên bố này cung cấp về diễn biến tương lai của chính sách tiền tệ.

Xét đến tầm quan trọng của việc chia sẻ dữ liệu giữa các nhà nghiên cứu và người thực hành, nhiều tổ chức đã bắt đầu nỗ lực hướng đến mục tiêu này. Ủy ban châu Âu (EC) đã khởi động nhiều sáng kiến, chẳng hạn như cổng thông tin EU Open Data\footnote{Cổng EU Open Data: \url{https://data.europa.eu/euodp/en/home/}.} và European Data\footnote{Cổng European Data: \url{https://www.europeandataportal.eu/en/homepage}.}, nhắm trực tiếp đến việc tạo điều kiện cho việc chia sẻ dữ liệu và khả năng tương tác.

\subsection{Quản lý dữ liệu và hạ tầng}

Để quản lý và phân tích khối lượng dữ liệu lớn xuất hiện ngày nay, cần có những hạ tầng mới có khả năng xử lý hiệu quả bốn chiều dữ liệu lớn: khối lượng, sự đa dạng, độ xác thực và tốc độ. Thật vậy, các tập dữ liệu khổng lồ cần được lưu trữ trong các môi trường điện toán phân tán chuyên biệt, vốn thiết yếu để xây dựng các đường ống dữ liệu (data pipeline) nhằm cắt lát và tổng hợp lượng thông tin lớn này. Dữ liệu phi cấu trúc lớn được lưu trữ trong các hệ thống tệp phân tán (Distributed File Systems — DFS), kết hợp nhiều máy tính (nút, node) thông qua một mạng lại với nhau~[36]. Dữ liệu được chia thành các khối (block) và lưu trữ trên các nút khác nhau, sao cho DFS cho phép làm việc với dữ liệu được phân mảnh, thứ mà nếu không sẽ trở nên quá lớn để lưu trữ và phân tích trên một máy tính đơn lẻ. Các khung công cụ sử dụng nhiều DFS bao gồm Apache Hadoop\footnote{Apache Hadoop: \url{https://hadoop.apache.org/}.} và Amazon S3\footnote{Amazon AWS S3: \url{https://aws.amazon.com/s3/}.}, nền tảng lưu trữ cốt lõi của AWS. Có rất nhiều nền tảng để xử lý và phân tích dữ liệu phân tán, trong đó nổi bật nhất có lẽ là Apache Spark\footnote{Apache Spark: \url{https://spark.apache.org/}.}. Khi làm việc với dữ liệu lớn, cần sử dụng các thuật toán chuyên biệt tránh việc phải nạp toàn bộ dữ liệu vào bộ nhớ làm việc của máy tính cùng một lúc~[36]. Chẳng hạn, khung công cụ MapReduce\footnote{\url{https://hadoop.apache.org/docs/r1.2.1/mapred_tutorial.html}.} bao gồm một chuỗi các thuật toán có thể chuẩn bị và nhóm dữ liệu thành các phần nhỏ tương đối (Map) trước khi thực hiện phân tích trên từng phần (Reduce). Các nền tảng DFS phổ biến khác hiện nay bao gồm MongoDB\footnote{MongoDB: \url{https://www.mongodb.com/}.}, Apache Cassandra\footnote{Apache Cassandra: \url{https://cassandra.apache.org/}.}, và ElasticSearch\footnote{ElasticSearch: \url{https://www.elastic.co/}.}, chỉ là một vài ví dụ tiêu biểu. Là một ví dụ trong lĩnh vực kinh tế, các tác giả của~[38] đã trình bày một hạ tầng NO-SQL dựa trên ElasticSearch để lưu trữ và tương tác với khối lượng khổng lồ dữ liệu tin tức chứa trong Cơ sở dữ liệu Toàn cầu về Sự kiện, Ngôn ngữ và Sắc thái (Global Database of Events, Language and Tone — GDELT)\footnote{Trang GDELT: \url{https://blog.gdeltproject.org/}.}, bao gồm hơn 8 TB thông tin văn bản từ khoảng 500 triệu bài báo tin tức trên toàn thế giới kể từ năm 2015. Các tác giả đã trình bày một ứng dụng khai thác GDELT để xây dựng các thước đo tâm lý tài chính dựa trên tin tức, nắm bắt quan điểm của nhà đầu tư đối với ba quốc gia châu Âu: Ý, Tây Ban Nha và Pháp~[38].

Mặc dù nhiều nền tảng dữ liệu lớn này cung cấp các giải pháp phù hợp cho doanh nghiệp và tổ chức để xử lý lượng dữ liệu và thông tin ngày càng tăng, nhiều ứng dụng quan trọng vẫn chưa được thiết kế để có khả năng mở rộng linh hoạt (dynamically scalable), cho phép tính toán phân tán, làm việc với các cơ sở dữ liệu phi truyền thống, hoặc tương tác với các hạ tầng khác. Các hạ tầng đám mây hiện có sẽ phải đầu tư mạnh mẽ vào các giải pháp được thiết kế để cung cấp khả năng mở rộng linh hoạt, khả năng tương tác giữa các hạ tầng, và tính toán song song quy mô lớn, nhằm cho phép thực thi đáng tin cậy các thuật toán học máy và kỹ thuật AI. Trong số các hành động khác, tầm quan trọng của điện toán đám mây gần đây đã được EC nhấn mạnh thông qua Sáng kiến Đám mây Châu Âu (European Cloud Initiative)\footnote{Sáng kiến Đám mây Châu Âu: \url{https://ec.europa.eu/digital-single-market/en/\%20european-cloud-initiative}.}, dẫn đến sự ra đời của Đám mây Khoa học Mở Châu Âu (European Open Science Cloud)\footnote{Đám mây Khoa học Mở Châu Âu: \url{https://ec.europa.eu/research/openscience/index.cfm?pg=open-science-cloud}.}, một môi trường mở đáng tin cậy dành cho cộng đồng khoa học để lưu trữ, chia sẻ và tái sử dụng dữ liệu và kết quả khoa học, cũng như Hạ tầng Dữ liệu Châu Âu (European Data Infrastructure)\footnote{Hạ tầng Dữ liệu Châu Âu: \url{https://www.eudat.eu/}.}, hướng đến việc xây dựng năng lực siêu máy tính của EU.

\section{Các phương pháp phân tích dữ liệu}

Các mô hình dự báo và ước tính hiện tại kinh tế truyền thống không có khả năng mở rộng linh hoạt để quản lý và duy trì các cấu trúc dữ liệu lớn, bao gồm nhật ký thô về hành động người dùng, văn bản tự nhiên từ giao tiếp, hình ảnh, video và dữ liệu cảm biến. Khối lượng dữ liệu lớn này xuất hiện dưới các định dạng nhiều chiều phức tạp về bản chất, và việc sử dụng chúng cho phân tích kinh tế đòi hỏi những bộ công cụ mới~[36]. Trên thực tế, các kỹ thuật truyền thống không mở rộng tốt khi số chiều dữ liệu lớn hoặc tăng nhanh. Những nhiệm vụ tương đối đơn giản như trực quan hóa dữ liệu, khớp mô hình (model fitting), và kiểm tra hiệu năng trở nên khó khăn. Các kiểm định giả thuyết cổ điển nhằm kiểm tra tầm quan trọng của một biến trong mô hình (kiểm định T), hoặc để chọn một mô hình trong số nhiều lựa chọn thay thế (kiểm định F), cần được sử dụng thận trọng trong môi trường dữ liệu lớn~[26, 30]. Trong bối cảnh phức tạp này, không thể dựa vào những đảm bảo chính xác đối với các chiến lược ít chiều tiêu chuẩn, các phương pháp trực quan hóa, và các chẩn đoán đặc tả mô hình~[36, 26]. Trong những bối cảnh này, các nhà khoa học xã hội có thể hưởng lợi từ việc sử dụng các kỹ thuật khoa học dữ liệu, và trong những năm gần đây, nỗ lực làm cho các ứng dụng đó được chấp nhận trong lĩnh vực mô hình hóa kinh tế đã tăng lên đáng kể. Một trọng tâm là mở "hộp đen" của các giải pháp học máy và xây dựng những mô hình có thể diễn giải được~[22]. Thật vậy, các thuật toán khoa học dữ liệu trở nên vô dụng đối với việc hoạch định chính sách khi, dù có khả năng mở rộng dễ dàng và hiệu năng cao, chúng lại khó có thể hiểu được. Khoa học dữ liệu tốt khi áp dụng vào kinh tế và tài chính đòi hỏi sự cân bằng giữa các khía cạnh này, và thường liên quan đến sự kết hợp giữa kiến thức chuyên môn và các công cụ phân tích nhằm đạt được mức độ hiệu năng mô hình, khả năng diễn giải, và mức độ tự động hóa mà các bên liên quan yêu cầu. Do đó, một thực hành tốt đối với các nhà kinh tế là xác định điều gì có thể được mô hình hóa như một nhiệm vụ dự báo, và dành các nỗ lực thống kê và kinh tế cho những câu hỏi mang tính cấu trúc khó khăn hơn. Trong phần tiếp theo, chúng tôi cung cấp một cái nhìn tổng quan ở mức cao về có lẽ là hai họ công nghệ khoa học dữ liệu phổ biến nhất được sử dụng ngày nay trong kinh tế và tài chính.

\subsection{Học máy sâu (Deep Machine Learning)}

Mặc dù các công nghệ học máy đã được thiết lập từ lâu, như Máy vector hỗ trợ (Support Vector Machines), Cây quyết định (Decision Trees), Rừng ngẫu nhiên (Random Forests), và Gradient Boosting đã cho thấy tiềm năng cao trong việc giải quyết nhiều bài toán khai phá dữ liệu (ví dụ, phân loại, hồi quy) xoay quanh các tổ chức, chính phủ và cá nhân. Ngày nay, công nghệ đạt được thành công lớn nhất trong cả giới nghiên cứu lẫn giới thực hành là học sâu (deep learning)~[29]. Học sâu là một công nghệ học máy đa năng, thường đề cập đến một tập hợp các thuật toán học máy dựa trên việc học các biểu diễn dữ liệu (data representation), nắm bắt các mối quan hệ phi tuyến bậc cao của dữ liệu đầu vào phi cấu trúc ở mức thấp để hình thành các khái niệm ở mức cao. Các phương pháp học sâu đã tạo ra bước đột phá thực sự về hiệu năng trong nhiều nhiệm vụ ở các lĩnh vực khác nhau mà các phương pháp học máy truyền thống gặp khó khăn, chẳng hạn như nhận dạng giọng nói, dịch máy, và thị giác máy tính (nhận dạng đối tượng). Ưu điểm của các thuật toán học sâu là khả năng phân tích dữ liệu rất phức tạp, chẳng hạn như hình ảnh, video, văn bản và các dữ liệu phi cấu trúc khác.

Các mô hình phân cấp sâu là các Mạng nơ-ron nhân tạo (Artificial Neural Networks — ANN) với cấu trúc sâu và các phương pháp liên quan, chẳng hạn như Deep Restricted Boltzmann Machines, Deep Belief Networks, và Mạng nơ-ron tích chập sâu (Deep Convolutional Neural Networks). ANN là các công cụ tính toán có thể được xem như lấy cảm hứng từ cách bộ não hoạt động, và áp dụng khung khái niệm này để xây dựng các mô hình toán học~[30]. Mạng nơ-ron ước lượng các hàm số có độ phức tạp tùy ý dựa trên dữ liệu đã cho. Mạng nơ-ron có giám sát được sử dụng để biểu diễn một ánh xạ từ một vector đầu vào sang một vector đầu ra. Mạng nơ-ron không giám sát thì được sử dụng để phân loại dữ liệu mà không cần kiến thức trước về các lớp liên quan. Về bản chất, Mạng nơ-ron có thể được xem như các mô hình hồi quy tổng quát hóa, có khả năng mô hình hóa dữ liệu với độ phức tạp tùy ý~[30]. Các kiến trúc ANN phổ biến nhất là perceptron đa lớp (multilayer perceptron — MLP) và hàm cơ sở bán kính (radial basis function — RBF). Trong thực tế, các chuỗi lớp ANN nối tiếp nhau tạo thành một khung học sâu. Thành công hiện tại của các phương pháp học sâu được thúc đẩy bởi những tiến bộ trong thuật toán và công nghệ điện toán hiệu năng cao, cho phép phân tích các tập dữ liệu lớn hiện đã có sẵn. Một ví dụ là các công cụ robot-advisor hiện đang sử dụng công nghệ học sâu để cải thiện độ chính xác~[19]. Các công cụ này thực hiện dự báo thị trường chứng khoán bằng cách giải quyết bài toán hồi quy hoặc bằng cách chuyển đổi thành bài toán phân loại và dự báo liệu thị trường sẽ tăng hay giảm.

Cũng có một khối lượng lớn tài liệu về việc sử dụng học sâu trong bối cảnh dự báo chuỗi thời gian~[29, 6, 27, 5]. Mặc dù việc sử dụng MLP ANN cổ điển trên các tập dữ liệu lớn khá đơn giản, việc sử dụng nó trên các chuỗi thời gian có kích thước trung bình lại khó khăn hơn do rủi ro cao về hiện tượng quá khớp (overfitting). MLP cổ điển có thể được điều chỉnh để xử lý bản chất tuần tự của dữ liệu bằng cách coi thời gian như một phần rõ ràng của đầu vào. Tuy nhiên, cách tiếp cận như vậy có một số khó khăn cố hữu, cụ thể là không có khả năng xử lý các chuỗi có độ dài thay đổi và phát hiện các mẫu bất biến theo thời gian trong dữ liệu. Một cách tiếp cận trực tiếp hơn là sử dụng các kết nối hồi quy (recurrent connections), kết nối các đơn vị ẩn của mạng nơ-ron trở lại chính chúng với một độ trễ thời gian. Đây là nguyên lý nền tảng của Mạng nơ-ron hồi quy (Recurrent Neural Networks — RNN)~[29], và đặc biệt là của Mạng bộ nhớ dài-ngắn hạn (Long Short-Term Memory Networks — LSTM)~[25], các ANN được thiết kế chuyên biệt để xử lý dữ liệu tuần tự xuất hiện trong các ứng dụng như chuỗi thời gian, xử lý ngôn ngữ tự nhiên, và nhận dạng giọng nói~[34].

Trong lĩnh vực tài chính, học sâu đã được khai thác, ví dụ, cho phân tích và dự báo thị trường chứng khoán (xem, ví dụ,~[13] để có một bài tổng quan). Một cách tiếp cận ANN đã được chứng minh khác cho dự báo chuỗi thời gian tài chính là Mạng nơ-ron tích chập giãn (Dilated Convolutional Neural Network) được trình bày trong~[9], trong đó kiến trúc nền tảng bắt nguồn từ dự án WaveNet của DeepMind~[41]. Công trình trong~[5] khai thác một tập hợp (ensemble) các Mạng nơ-ron tích chập, được huấn luyện trên các ảnh Gramian Angular Fields được sinh ra từ chuỗi thời gian liên quan đến chỉ số Standard \& Poor's 500 Future, với mục tiêu là dự báo xu hướng tương lai của thị trường Mỹ.

Bên cạnh học sâu, \emph{học tăng cường} (reinforcement learning) đã trở nên phổ biến trong những năm gần đây: đây là một mô hình học dựa trên thử và sai, chỉ dựa vào phần thưởng hoặc hình phạt. Nó đã được áp dụng thành công trong những đổi mới mang tính đột phá, chẳng hạn như hệ thống AlphaGo\footnote{Hệ thống AlphaGo của Deep Mind: \url{https://deepmind.com/research/case-studies/alphago-the-story-so-far}.} của Deep Mind đã chiến thắng người chơi cờ vây giỏi nhất thế giới. Nó cũng có thể được áp dụng trong lĩnh vực kinh tế, ví dụ để tối ưu hóa danh mục đầu tư một cách linh hoạt~[23] hoặc cho giao dịch tài sản tài chính~[18]. Tất cả những hệ thống học máy tiên tiến này có thể được sử dụng để học và liên kết thông tin từ nhiều nguồn kinh tế khác nhau, và xác định những mối tương quan ẩn không thể nhìn thấy khi chỉ xem xét một nguồn dữ liệu duy nhất. Chẳng hạn, việc kết hợp các đặc trưng từ hình ảnh (ví dụ, ảnh vệ tinh) và văn bản (ví dụ, mạng xã hội) có thể mang lại sự cải thiện cho dự báo kinh tế.

Việc xây dựng một quy trình học sâu hoặc học tăng cường hoàn chỉnh, bao gồm các nhiệm vụ quan trọng như xử lý dữ liệu, diễn giải, thiết kế khung công cụ, và tinh chỉnh tham số, thiên về nghệ thuật (hoặc một kỹ năng học được từ kinh nghiệm) hơn là một khoa học chính xác. Tuy nhiên, công việc này được hỗ trợ bởi các ngôn ngữ lập trình được sử dụng để phát triển những quy trình như vậy, ví dụ R, Scala và Python, cung cấp không gian làm việc tuyệt vời cho nhiều ứng dụng khoa học dữ liệu, đặc biệt là những ứng dụng liên quan đến dữ liệu phi cấu trúc. Những ngôn ngữ lập trình này đang tiến lên các cấp độ cao hơn, nghĩa là giờ đây có thể sử dụng những chỉ lệnh ngắn gọn và trực quan để tự động giải quyết một số vấn đề lập trình phức tạp và tốn công sức, ví dụ như cấp phát bộ nhớ, phân mảnh dữ liệu, và tối ưu hóa tham số. Chẳng hạn, thư viện Gluon\footnote{Gluon: \url{https://gluon.mxnet.io/}.} hiện đang phổ biến, bọc (tức là cung cấp chức năng cấp cao hơn xung quanh) MXNet\footnote{Apache MXNet: \url{https://mxnet.apache.org/}.}, một khung học sâu giúp việc xây dựng mạng nơ-ron sâu trở nên dễ dàng và nhanh chóng hơn. Bản thân MXNet bọc C++, mã nguồn nhanh và hiệu quả về bộ nhớ, được biên dịch thực sự để thực thi. Tương tự, Keras\footnote{Keras: \url{https://keras.io/}.}, một thư viện được sử dụng rộng rãi khác, là một phần mở rộng của Python bọc lại một số khung học sâu khác, chẳng hạn như TensorFlow\footnote{TensorFlow: \url{https://www.tensorflow.org/}.} của Google. Những công cụ này và các công cụ trong tương lai đang tạo ra một thế giới với các giao diện thân thiện với người dùng để học máy (sâu) nhanh hơn và đơn giản hơn~[36].

\subsection{Công nghệ Web ngữ nghĩa (Semantic Web Technologies)}

Từ góc độ xử lý và khai phá nội dung dữ liệu, dữ liệu văn bản thuộc về cái gọi là dữ liệu phi cấu trúc. Việc học từ loại dữ liệu phức tạp này có thể tạo ra những mẫu hình mô tả cô đọng và giàu ngữ nghĩa hơn, phản ánh tốt hơn các thuộc tính nội tại của dữ liệu. Các công nghệ như Web ngữ nghĩa, bao gồm Xử lý ngôn ngữ tự nhiên (Natural Language Processing — NLP) và Tìm kiếm thông tin (Information Retrieval), đã được tạo ra để tạo điều kiện dễ dàng tiếp cận nguồn thông tin văn bản phong phú. Web ngữ nghĩa, thường được gọi là "Web 3.0", là một hệ thống cho phép máy móc "hiểu" và phản hồi các yêu cầu phức tạp của con người dựa trên ý nghĩa của chúng. Sự "hiểu biết" như vậy đòi hỏi các nguồn thông tin liên quan phải được cấu trúc về mặt ngữ nghĩa~[7]. Linked Open Data (LOD) đã có được động lực đáng kể trong những năm qua như một thực hành tốt nhất để thúc đẩy việc chia sẻ và công bố dữ liệu có cấu trúc trên Web ngữ nghĩa~[8], bằng cách cung cấp một mô tả chính thức về các khái niệm, thuật ngữ và mối quan hệ trong một lĩnh vực tri thức nhất định, và bằng cách sử dụng Định danh tài nguyên thống nhất (Uniform Resource Identifiers — URI), Khung mô tả tài nguyên (Resource Description Framework — RDF), và Ngôn ngữ bản thể Web (Web Ontology Language — OWL), các tiêu chuẩn này đều do W3C quản lý.

LOD mang lại khả năng sử dụng dữ liệu giữa các lĩnh vực khác nhau cho các mục đích như thống kê, phân tích, bản đồ và xuất bản. Bằng cách liên kết tri thức này, các mối liên hệ và sự kết hợp có thể được suy ra và những kết luận mới có thể được rút ra. RDF/OWL cho phép tạo ra các bộ ba (triples) về bất kỳ điều gì trên Web ngữ nghĩa: không gian dữ liệu phi tập trung của tất cả các bộ ba này đang tăng trưởng với tốc độ đáng kinh ngạc, vì ngày càng có nhiều nguồn dữ liệu được công bố dưới dạng dữ liệu ngữ nghĩa. Nhưng quy mô của Web ngữ nghĩa không phải là tham số duy nhất của độ phức tạp ngày càng tăng của nó. Đặc tính phân tán và động của nó, cùng với các vấn đề về tính nhất quán giữa các nguồn dữ liệu, và sự tương tác giữa các nguồn dữ liệu thông qua suy luận (reasoning), góp phần biến Web ngữ nghĩa thành một hệ thống lớn, phức tạp~[7, 8].

Một trong những công nghệ phổ biến nhất được dùng để giải quyết các nhiệm vụ khác nhau trong Web ngữ nghĩa là NLP, thường được gọi bằng các từ đồng nghĩa như khai phá văn bản (text mining), phân tích văn bản (text analytics), hay khám phá tri thức từ văn bản (knowledge discovery from text). NLP là một thuật ngữ rộng, đề cập đến các công nghệ và phương pháp trong ngôn ngữ học tính toán nhằm tự động phát hiện và phân tích thông tin liên quan trong nội dung văn bản phi cấu trúc (văn bản tự do). Đã có những bước đột phá đáng kể trong NLP với sự ra đời của các công nghệ học máy tiên tiến (đặc biệt là học sâu) và các phương pháp thống kê cho các nhiệm vụ phân tích văn bản chính, như: phân tích ngôn ngữ học, nhận dạng thực thể có tên (named entity recognition), giải quyết đồng tham chiếu (co-reference resolution), trích xuất quan hệ (relation extraction), và phân tích quan điểm và cảm xúc (opinion and sentiment analysis)~[15].

Trong lĩnh vực kinh tế, các công cụ NLP đã được điều chỉnh và phát triển thêm để trích xuất các khái niệm, cảm xúc và tâm trạng liên quan từ mạng xã hội và tin tức (xem, ví dụ,~[37, 24, 14, 4], trong số nhiều tài liệu khác). Những công nghệ này khi được áp dụng trong bối cảnh kinh tế tạo điều kiện thuận lợi cho việc tích hợp dữ liệu từ nhiều nguồn không đồng nhất, cho phép phát triển các hệ thống lọc thông tin, và hỗ trợ các nhiệm vụ khám phá tri thức.

\section{Kết luận}

Trong chương này, chúng tôi đã giới thiệu chủ đề khoa học dữ liệu ứng dụng vào mô hình hóa kinh tế và tài chính. Các thách thức như xử lý dữ liệu kinh tế, chất lượng, số lượng, bảo vệ, và tích hợp đã được trình bày, cùng với các hạ tầng quản lý dữ liệu lớn chính yếu và các phương pháp phân tích dữ liệu phục vụ cho dự báo, diễn giải, khai phá, và các nhiệm vụ khám phá tri thức. Chúng tôi đã tóm tắt một số vấn đề phổ biến về dữ liệu lớn trong mô hình hóa kinh tế và các phương pháp khoa học dữ liệu liên quan.

Rõ ràng có nhu cầu và tiềm năng cao để phát triển các phương pháp khoa học dữ liệu cho phép con người và máy móc hợp tác chặt chẽ hơn, nhằm đạt được những mô hình được cải thiện trong kinh tế và tài chính. Những công nghệ này có thể xử lý, phân tích và khai thác tập hợp dữ liệu rất đa dạng, liên kết chặt chẽ và phức tạp vốn đã tồn tại trong vũ trụ kinh tế, nhằm cải thiện các mô hình và chất lượng dự báo, về mặt đảm bảo độ tin cậy của thông tin, tập trung vào việc tạo ra các khuyến nghị có thể hành động (actionable advice), và cải thiện tính tương tác của việc xử lý và phân tích dữ liệu.

\section*{Tài liệu tham khảo}

\begin{enumerate}
\item Aruoba, S. B., Diebold, F. X., \& Scotti, C. (2009). Real-time measurement of business conditions. \emph{Journal of Business \& Economic Statistics}, 27(4), 417--427.
\item Babii, A., Chen, X., \& Ghysels, E. (2019). Commercial and residential mortgage defaults: Spatial dependence with frailty. \emph{Journal of Econometrics}, 212, 47--77.
\item Baesens, B., Van Vlasselaer, V., \& Verbeke, W. (2015). \emph{Fraud analytics using descriptive, predictive, and social network techniques: a guide to data science for fraud detection}. Chichester: John Wiley \& Sons.
\item Barbaglia, L., Consoli, S., \& Manzan, S. (2020). Monitoring the business cycle with fine-grained, aspect-based sentiment extraction from news. Trong V. Bitetta et al. (chủ biên), \emph{Mining Data for Financial Applications (MIDAS 2019)}, Lecture Notes in Computer Science (Vol. 11985, tr. 101--106). Cham: Springer.
\item Barra, S., Carta, S., Corriga, A., Podda, A. S., \& Reforgiato Recupero, D. (2020). Deep learning and time series-to-image encoding for financial forecasting. \emph{IEEE Journal of Automatica Sinica}, 7, 683--692.
\item Benidis, K., Rangapuram, S. S., Flunkert, V., Wang, B., Maddix, D. C., Türkmen, C., Gasthaus, J., Bohlke-Schneider, M., Salinas, D., Stella, L., Callot, L., \& Januschowski, T. (2020). Neural forecasting: Introduction and literature overview. \emph{CoRR}, abs/2004.10240.
\item Berners-Lee, T., Chen, Y., Chilton, L., Connolly, D., Dhanaraj, R., Hollenbach, J., Lerer, A., \& Sheets, D. (2006). Tabulator: Exploring and analyzing linked data on the semantic web. Trong \emph{Proc. 3rd International Semantic Web User Interaction Workshop (SWUI 2006)}.
\item Bizer, C., Heath, T., \& Berners-Lee, T. (2009). Linked Data - The story so far. \emph{International Journal on Semantic Web and Information Systems}, 5, 1--22.
\item Borovykh, A., Bohte, S., \& Oosterlee, C. W. (2017). Conditional time series forecasting with convolutional neural networks. \emph{Lecture Notes in Computer Science}, 10614, 729--730.
\item Buneman, P., \& Tan, W.-C. (2019). Data provenance: What next? \emph{ACM SIGMOD Record}, 47(3), 5--16.
\item Carta, S., Fenu, G., Reforgiato Recupero, D., \& Saia, R. (2019). Fraud detection for e-commerce transactions by employing a prudential multiple consensus model. \emph{Journal of Information Security and Applications}, 46, 13--22.
\item Carta, S., Consoli, S., Piras, L., Podda, A. S., \& Reforgiato Recupero, D. (2020). Dynamic industry specific lexicon generation for stock market forecast. Trong G. Nicosia et al. (chủ biên), \emph{Machine Learning, Optimization, and Data Science (LOD 2020)}, Lecture Notes in Computer Science (Vol. 12565, tr. 162--176). Cham: Springer.
\item Chong, E., Han, C., \& Park, F. C. (2017). Deep learning networks for stock market analysis and prediction: Methodology, data representations, and case studies. \emph{Expert Systems with Applications}, 83, 187--205.
\item Consoli, S., Tiozzo Pezzoli, L., \& Tosetti, E. (2020). Using the GDELT dataset to analyse the Italian bond market. Trong G. Nicosia et al. (chủ biên), \emph{Machine learning, optimization, and data science (LOD 2020)}, Lecture Notes in Computer Science (Vol. 12565, tr. 190--202). Cham: Springer.
\item Consoli, S., Reforgiato Recupero, D., \& Petkovic, M. (2019). \emph{Data science for healthcare - Methodologies and applications}. Berlin: Springer Nature.
\item Daily, J., \& Peterson, J. (2017). Predictive maintenance: How big data analysis can improve maintenance. Trong \emph{Supply chain integration challenges in commercial aerospace} (tr. 267--278). Cham: Springer.
\item Dal Pozzolo, A., Caelen, O., Johnson, R. A., \& Bontempi, G. (2015). Calibrating probability with undersampling for unbalanced classification. Trong \emph{2015 IEEE Symposium Series on Computational Intelligence} (tr. 159--166). Piscataway: IEEE.
\item Deng, Y., Bao, F., Kong, Y., Ren, Z., \& Dai, Q. (2017). Deep direct reinforcement learning for financial signal representation and trading. \emph{IEEE Transactions on Neural Networks and Learning Systems}, 28(3), 653--664.
\item Ding, X., Zhang, Y., Liu, T., \& Duan, J. (2015). Deep learning for event-driven stock prediction. Trong \emph{IJCAI International Joint Conference on Artificial Intelligence} (Vol. 2015, tr. 2327--2333).
\item Ertan, A., Loumioti, M., \& Wittenberg-Moerman, R. (2017). Enhancing loan quality through transparency: Evidence from the European central bank loan level reporting initiative. \emph{Journal of Accounting Research}, 55(4), 877--918.
\item Giannone, D., Reichlin, L., \& Small, D. (2008). Nowcasting: The real-time informational content of macroeconomic data. \emph{Journal of Monetary Economics}, 55(4), 665--676.
\item Gilpin, L. H., Bau, D., Yuan, B. Z., Bajwa, A., Specter, M., \& Kagal, L. (2019). Explaining explanations: An overview of interpretability of machine learning. Trong \emph{IEEE International Conference on Data Science and Advanced Analytics (DSAA 2018)} (tr. 80--89).
\item Goodfellow, I., Bengio, Y., \& Courville, A. (2016). \emph{Deep Learning}. Cambridge: MIT Press.
\item Hansen, S., \& McMahon, M. (2016). Shocking language: Understanding the macroeconomic effects of central bank communication. \emph{Journal of International Economics}, 99, S114--S133.
\item Hochreiter, S., \& Schmidhuber, J. (1997). Long short-term memory. \emph{Neural Computation}, 9, 1735--1780.
\item Jabbour, C. J. C., Jabbour, A. B. L. D. S., Sarkis, J., \& Filho, M. G. (2019). Unlocking the circular economy through new business models based on large-scale data: An integrative framework and research agenda. \emph{Technological Forecasting and Social Change}, 144, 546--552.
\item Januschowski, T., Gasthaus, J., Wang, Y., Salinas, D., Flunkert, V., Bohlke-Schneider, M., \& Callot, L. (2020). Criteria for classifying forecasting methods. \emph{International Journal of Forecasting}, 36(1), 167--177.
\item Kuzin, V., Marcellino, M., \& Schumacher, C. (2011). MIDAS vs. mixed-frequency VAR: Nowcasting GDP in the euro area. \emph{International Journal of Forecasting}, 27(2), 529--542.
\item LeCun, Y., Bengio, Y., \& Hinton, G. (2015). Deep Learning. \emph{Nature}, 521(7553), 436--444.
\item Marwala, T. (2013). \emph{Economic modeling using Artificial Intelligence methods}. Heidelberg: Springer.
\item Marx, V. (2013). The big challenges of big data. \emph{Nature}, 498, 255--260.
\item Oblé, F., \& Bontempi, G. (2019). Deep-learning domain adaptation techniques for credit cards fraud detection. Trong \emph{Recent Advances in Big Data and Deep Learning: Proceedings of the INNS Big Data and Deep Learning Conference} (Vol. 1, tr. 78--88). Cham: Springer.
\item OECD. (2015). \emph{Data-driven innovation: Big data for growth and well-being}. OECD Publishing, Paris.
\item Salinas, D., Flunkert, V., Gasthaus, J., \& Januschowski, T. (2020). Deepar: Probabilistic forecasting with autoregressive recurrent networks. \emph{International Journal of Forecasting}, 36(3), 1181--1191.
\item Sirignano, J., Sadhwani, A., \& Giesecke, K. (2018). \emph{Deep learning for mortgage risk}. Technical report, Working paper available at SSRN.
\item Taddy, M. (2019). \emph{Business data science: Combining machine learning and economics to optimize, automate, and accelerate business decisions}. New York: McGraw-Hill, US.
\item Tetlock, P. C. (2007). Giving content to investor sentiment: The role of media in the stock market. \emph{The Journal of Finance}, 62(3), 1139--1168.
\item Tiozzo Pezzoli, L., Consoli, S., \& Tosetti, E. (2020). Big data financial sentiment analysis in the European bond markets. Trong V. Bitetta et al. (chủ biên), \emph{Mining Data for Financial Applications (MIDAS 2019)}, Lecture Notes in Computer Science (Vol. 11985, tr. 122--126). Cham: Springer.
\item Tiwari, S., Wee, H. M., \& Daryanto, Y. (2018). Big data analytics in supply chain management between 2010 and 2016: Insights to industries. \emph{Computers \& Industrial Engineering}, 115, 319--330.
\item Van Bekkum, S., Gabarro, M., \& Irani, R. M. (2017). Does a larger menu increase appetite? Collateral eligibility and credit supply. \emph{The Review of Financial Studies}, 31(3), 943--979.
\item van den Oord, A., Dieleman, S., Zen, H., Simonyan, K., Vinyals, O., Graves, A., et al. (2016). WaveNet: A generative model for raw audio. \emph{CoRR}, abs/1609.03499.
\item Wilkinson, M., Dumontier, M., Aalbersberg, I., Appleton, G., Axton, M., Baak, A., et al. (2016). The FAIR guiding principles for scientific data management and stewardship. \emph{Scientific Data}, 3, 1.
\item Wu, X., Zhu, X., Wu, G., \& Ding, W. (2014). Data mining with Big Data. \emph{IEEE Transactions on Knowledge and Data Engineering}, 26(1), 97--107.
\end{enumerate}
